\part{Ordered collisions and reconstruction}

\chapter{Ordered words and collision coefficients}
\label{chap:collision-words}

Three ordered insertions can have the same first and third position while
the second remains separate.  A construction that remembers only adjacent
collisions loses this configuration.  A construction that remembers only
the three singleton state spaces loses the normal derivatives along its
collision diagonal.  Both kinds of information enter reconstruction.

Let \(k\) be a field of characteristic zero and let \(X\) be a smooth
separated curve over \(k\), and write \([n]=\{1,\ldots,n\}\).
The positions of a word of length \(n\) form \(X^n\), including its
diagonals.  The order belongs to the labels in \([n]\).
It does not order the points of \(X\).  Cutting the labels
between \(p\) and \(p+1\) is permitted where every point of the first
block is distinct from every point of the second.  Collisions inside
either block remain permitted.

\section{The coefficient category}

We use cochain complexes: the differential has degree \(1\), and
\(M[r]^i=M^{i+r}\).  Thus an element \(sa\) of \(M[1]\) has degree
\(|a|-1\).  Tensor products of maps and differentials use the Koszul
sign rule.  All tensor products of complexes are derived unless a flat
model has been specified.

Fix dg-enhanced stable coefficient categories \(\mathscr D(X^n)\), with
external tensor products, exceptional inverse images \(f^!\), and
direct images \(j_*\) for open embeddings.  These functors are exact,
and external tensor products preserve colimits separately.
We require the following
six-functor identities on the diagrams used below: open base change,
composition of open direct images, compatibility of external products
with those operations, and localization for a closed complement.  For
the diagonals that identify consecutive labels we also require the
base-change identifications in \eqref{eq:words-diagonal-cut}.  These
identities are part of the coefficient setting, with their coherent
enhancements.  They are the hypotheses on a differential-module or
sheaf realization of the constructions.

A monotone surjection \(\alpha:[n]\twoheadrightarrow[r]\) has interval
fibres.  Its diagonal map is
\[
 d_\alpha:X^r\longrightarrow X^n,
 \qquad (y_1,\ldots,y_r)\longmapsto
 (y_{\alpha(1)},\ldots,y_{\alpha(n)}).
\]
An \emph{ordered word coefficient} is a family
\(F_n\in\mathscr D(X^n)\), \(n\geq1\), with coherent equivalences
\(d_\alpha^!F_n\simeq F_r\).  We adjoin a separate coefficient
\(F_0\in D(k)\) for the empty word, where \(D(k)\) is the derived
category of \(k\)-complexes.  There is no collision map from a
nonempty word to that empty word.  Write \(\mathscr W_X\) for this
stable category of diagrams.

Stability concerns the ambient diagram category.  A factorization
condition selects objects whose specified structure maps are
equivalences, and need not define a stable subcategory.  For example,
suppose a coefficient has rank one on every separated configuration
and is strongly factorizing.  Its direct sum with itself has rank two
on a separated pair.  The external square of its singleton coefficient
has rank four.  The componentwise direct sum therefore cannot retain
the same strong factorization condition.

An arbitrary partition of \([n]\) still defines a diagonal in \(X^n\).
The restriction of \(F_n\) to it, its supported sections, and its normal
derivatives remain meaningful.  The definition of \(\mathscr W_X\)
does not identify such a restriction with a shorter word unless the
partition is consecutive.  Comparisons on the full collision diagram
must supply their own maps and coherences.

\section{Consecutive cuts}

For \(0\leq p\leq n\), set
\[
 U_p^{(n)}=
 \{(x_1,\ldots,x_n):x_i\neq x_j\text{ whenever }i\leq p<j\},
 \qquad j_p:U_p^{(n)}\hookrightarrow X^n.
\]
The two end cuts have \(U_0^{(n)}=U_n^{(n)}=X^n\).  Define
\begin{equation}
 (F\star_c G)_n=
 \bigoplus_{p=0}^n
 j_{p,*}j_p^!(F_p\boxtimes G_{n-p}).
 \label{eq:words-cut-product}
\end{equation}
The coefficient \(K\), equal to \(k\) at the empty word and zero at
every nonempty word, will be its unit.

Here is the diagonal compatibility required of the coefficient
setting.  If \(p\) lies strictly inside a fibre of \(\alpha\), the
image of \(d_\alpha\) misses \(U_p^{(n)}\).  Exceptional pullback of
the corresponding localized term is zero.  Otherwise
\(p=|\alpha^{-1}([q])|\) for a unique cut \(q\) of \([r]\), and
the surviving base-change map gives
\begin{equation}
 d_\alpha^!j_{p,*}j_p^!(F_p\boxtimes G_{n-p})
 \simeq j_{q,*}j_q^!(F_q\boxtimes G_{r-q}).
 \label{eq:words-diagonal-cut}
\end{equation}
The indices on the right use the word-coefficient identifications.
We require these maps to commute with composition of monotone
surjections.  Summing them makes \eqref{eq:words-cut-product} an
object of \(\mathscr W_X\).

\begin{proposition}[The cut product]
\label{prop:words-cut-monoidal}
Under the stated coefficient identities, \((\mathscr W_X,\star_c,K)\)
is a monoidal stable category.  The product is exact separately in
each variable.  No interchange of the two ordered blocks is part of
this monoidal structure.
\end{proposition}
\begin{proof}
An iterated product of \(r\) factors is indexed by a consecutive
decomposition of the word into \(r\) blocks, with empty blocks allowed.
For any parenthesization, its coefficient is the external product of
the block coefficients, localized to the open where distinct blocks
have disjoint positions.  Open base change identifies these
coefficients.  For four or more factors, every composite associator
is the same identification with the coefficient of the common block
decomposition.  This proves the pentagon and, in the enhanced
categories, the higher coherence identities.  An empty factor
contributes only the corresponding end cut, giving both unit maps.
Exactness follows from the assumed exact operations and the finite
sum in each arity.
\end{proof}

Call \(F\) \emph{positive} if \(F_0=0\).  Positivity bounds the
number of factors that can contribute at a fixed word length.

\begin{lemma}[Arity bound]
\label{lem:words-arity-bound}
If \(F^1,\ldots,F^r\) are positive, then
\[
 (F^1\star_c\cdots\star_c F^r)_n=0
 \qquad\text{for }r>n.
\]
\end{lemma}
\begin{proof}
A nonzero summand requires a nonempty block for every factor.
A word of length \(n\) has at most \(n\) such blocks.
\end{proof}

This bound does not assert a bound on the pole order or on normal
derivatives.  Each surviving coefficient may contain all of them.
It bounds the number of positive factors, which is the number needed
in the bar--cobar argument.

\section{Finite normalized bar and cobar}

An augmented associative algebra in \(\mathscr W_X\) has unit and
augmentation \(K\to A\to K\).  Assume that \(A_0\to k\) is an
equivalence.  Its augmentation ideal \(\overline A\), the fibre of
\(A\to K\), is then positive.  For a coaugmented counital
coassociative coalgebra \(C\) with \(C_0=k\), write
\(\overline C=\operatorname{cofib}(K\to C)\).  The counit splits the
coaugmentation, so this is the reduced coefficient.

The normalized constructions have underlying coefficients
\begin{align}
 B_c A&=K\oplus\bigoplus_{r\geq1}
              (\overline A[1])^{\star_c r},
              \label{eq:words-bar}\\
 \Omega_c C&=K\oplus\bigoplus_{r\geq1}
              (\overline C[-1])^{\star_c r}.
              \label{eq:words-cobar}
\end{align}
Both sums are finite in each arity.  The first has deconcatenation
coproduct.  The second has concatenation product.  A flat differential
graded model makes their signs explicit.  Writing \(s\) for shift
by \([1]\), the bar differential is the coderivation with terms
\begin{equation}
 b_1(sa)=-s(da),\qquad
 b_2(sa\otimes sb)=(-1)^{|a|}s(ab).
 \label{eq:words-bar-signs}
\end{equation}
Its extension past earlier shifted factors uses their total degree.
If \(t\) denotes shift by \([-1]\) and
\(\overline\Delta c=\sum c'\otimes c''\), the cobar differential
is the derivation determined by
\begin{equation}
 d(tc)=-t(dc)+
       \sum(-1)^{|c'|+1}(tc')(tc'').
 \label{eq:words-cobar-signs}
\end{equation}
The derivation and coderivation identities, together with
associativity and coassociativity, give square-zero differentials.
For instance, the two contractions of three suspended factors in the
bar have signs \((-1)^{|b|}\) and \((-1)^{|b|-1}\).

\begin{theorem}[Bar--cobar on positive word coefficients]
\label{thm:words-finite-duality}
In the enhanced coefficient setting above, normalized bar and cobar
give inverse equivalences between augmented associative algebras
with empty coefficient \(k\) and coaugmented counital coassociative
coalgebras with empty coefficient \(k\).  The unit and counit are
equivalences in every arity.  No extra completion in bar length is
required for this statement.
\end{theorem}
\begin{proof}
Consider \(\Omega_c B_c A\) in arity \(n\).  Flatten each tensor of
bar words into its original positive factors from \(\overline A\).
There are at most \(n\) factors.  Filter by their number.  The
associated graded differential retains the internal differential
and the part that splits a bar word into two cobar factors.  It
omits the multiplication in \(A\), which reduces the number of
original factors.

For \(r\) original factors, each of the \(r-1\) gaps either lies
inside a bar word or separates two cobar factors.  Changing the
first state to the second is the identity of the same localized
coefficient, with the suspension sign.  Hence the gap complex is
the tensor product of \(r-1\) two-term identity complexes.  It is
contractible for \(r>1\): contract one chosen gap and leave the
others unchanged, using the tensor sign rule.  For \(r=1\) it is
\(\overline A\).  The counit sends \(t(sa)\) to \(a\), sends longer
single bar words to zero, and extends multiplicatively.  Formulae
\eqref{eq:words-bar-signs}--\eqref{eq:words-cobar-signs} make it a
chain map.  On the associated graded it is the identity on the
only noncontractible term.  The filtration is finite in arity
\(n\), so the counit is an equivalence there.

For \(B_c\Omega_c C\), use the number of original reduced
coalgebra factors with the descending filtration.  The reduced
coproduct increases this number, and concatenation preserves it.
The associated graded again consists of the same gap complexes,
with the direction of the gap differential reversed.  The unit on a
reduced coefficient has the formula
\[
 \eta(c)=\sum_{r\geq1}\sum
 [tc_{(1)}|\cdots|tc_{(r)}],
 \qquad
 \overline\Delta^{(r)}c=
       \sum c_{(1)}\otimes\cdots\otimes c_{(r)}.
\]
The brackets suspend each cobar generator once, so the composite
map \(c\mapsto s(tc)\) has degree zero.  The chain-map identity
follows from coassociativity and \eqref{eq:words-cobar-signs}.
Reduced iterates with more than \(n\) positive factors vanish by
Lemma~\ref{lem:words-arity-bound}.  The one-factor term
of the unit is the identity on \(\overline C\).  The same finite
filtration argument proves that the unit is an equivalence.  These
are the natural unit and counit of the bar--cobar adjunction, so
their triangle identities give the asserted inverse equivalences.
\end{proof}

An ordinary augmented algebra concentrated at a single position
does not satisfy this arity argument merely because it is augmented.
Its powers can survive in every tensor length.  Positive weight,
conilpotence or an appropriate completeness hypothesis must then
replace Lemma~\ref{lem:words-arity-bound}.

\section{The support left by all cuts}

A coalgebra \(C\) is \emph{strong} if, for every cut, its specified
coproduct restricts to an equivalence
\begin{equation}
 C_n|_{U_p^{(n)}}\longrightarrow
 (C_p\boxtimes C_{n-p})|_{U_p^{(n)}}.
 \label{eq:words-strong}
\end{equation}
These maps retain the coalgebra coherences.  A collection of
unrelated equivalences is insufficient.  Define the closed subset
\begin{equation}
 Z_n=X^n\setminus\bigcup_{p=1}^{n-1}U_p^{(n)}
     =\bigcap_{p=1}^{n-1}\ \bigcup_{i\leq p<j}\Delta_{ij},
 \qquad Z_1=X.
 \label{eq:words-support}
\end{equation}
Thus a configuration lies in \(Z_n\) precisely when every
consecutive cut is crossed by a coincident pair.

For a closed subset \(Z\subseteq Y\) with open complement
\(j:U\hookrightarrow Y\), the supported coefficient
\(R\Gamma_Z M\) is the fibre of \(M\to j_*j^!M\).

\begin{theorem}[Cut-supported reconstruction]
\label{thm:words-support}
Let \(C\) be a strong coalgebra with empty coefficient \(k\).
Then there is a natural equivalence of coefficients
\begin{equation}
 (\overline{\Omega_c C})_n
       \simeq R\Gamma_{Z_n}(C_n)[-1].
 \label{eq:words-supported-cobar}
\end{equation}
Conversely, if the augmentation ideal of \(A\) is supported on
\(Z_n\) in every arity, then \(B_cA\) is strong.  Consequently
bar--cobar restricts to an equivalence between strong coalgebras
and augmented algebras with these support conditions.
\end{theorem}
\begin{proof}
For a set of cuts \(S\subseteq\{1,\ldots,n-1\}\), put
\(U_S=\bigcap_{p\in S}U_p^{(n)}\), with \(U_\varnothing=X^n\),
and write \(L_S M=j_{S,*}j_S^!M\).  An \(r\)-factor term in
\(\overline{\Omega_c C}\) corresponds to \(r-1\) cuts.  Iterated
strong coproducts identify its coefficient with \(L_S C_n\).
Coassociativity identifies the splitting differential with the
maps that insert one more cut.  The resulting complex is
\[
 \operatorname{Tot}
 \left(\bigoplus_{S\subseteq[n-1]}L_SC_n[-|S|-1]\right),
\]
with the tensor suspension signs.  Apart from the common shift
\([-1]\), it is the total fibre of the localization cube for the
opens \(U_p^{(n)}\).  Finite open descent identifies that fibre
with \(R\Gamma_{Z_n}C_n\), proving
\eqref{eq:words-supported-cobar}.

For the converse, restrict a normalized bar word to \(U_p^{(n)}\).
If one of its consecutive factors crosses the cut \(p\), the
positions in that factor lie in its own cut open.  The support
hypothesis makes that factor vanish.  Every surviving word has
a separator at \(p\).  Cutting there is a bijection between
surviving bar words and pairs of bar words on the two blocks.
The bar differential has no remaining term that multiplies
across this separator, again by the support hypothesis.  Thus
deconcatenation is the equivalence \eqref{eq:words-strong},
including its differential and iterated-cut coherences.
The restricted equivalence follows from
Theorem~\ref{thm:words-finite-duality}.
\end{proof}

The right side of \eqref{eq:words-supported-cobar} describes the
coefficient.  Its multiplication still uses the inverses of the
strong coproducts and their coherence.  Local cohomology alone
does not specify that multiplication.

\begin{example}[Three and four positions]
\label{ex:words-crossed-support}
For three labels,
\[
 Z_3=(\Delta_{12}\cup\Delta_{13})
       \cap(\Delta_{13}\cup\Delta_{23})=\Delta_{13}.
\]
The locus with \(x_1=x_3\neq x_2\) is therefore part of the
reconstructed support.  For four labels,
\[
 Z_4=\Delta_{14}\ \cup\ (\Delta_{13}\cap\Delta_{24}).
\]
To check this equality, suppose \(x_1\neq x_4\).  The first cut
requires \(x_1=x_2\) or \(x_1=x_3\).  If \(x_1=x_2\), the last cut
forces \(x_3=x_4\), leaving the middle cut uncrossed.
Hence \(x_1=x_3\), and the last cut then forces \(x_2=x_4\).
On the affine line,
take the independent normal coordinates
\(a=x_1-x_4\), \(b=x_1-x_3\), \(c=x_2-x_4\).  The reduced ideal
of this union is
\[
 (a)\cap(b,c)=(ab,ac).
\]
Its codimension-one and codimension-two parts explain why a
single total-diagonal coefficient cannot replace the full support.
\end{example}

\section{The empty unit and the local vacuum}

Suppose the coefficient setting has dualizing objects
\(J_n=p_n^!k\), where \(p_n:X^n\to\operatorname{Spec}k\), with
the coherent product identifications
\(J_p\boxtimes J_q\simeq J_{p+q}\).  Set \(J_0=k\).  Localization
gives cut coproducts on \(J\).  They are equivalences on their cut
opens.  Thus \(J\) is a strong coalgebra.  Its cobar algebra
\(\mathcal U_X=\Omega_cJ\) carries the local vacuum coefficient.

A \emph{vacuum-marked coalgebra} has a coalgebra map \(J\to C\).
Applying bar--cobar gives the equivalence of undercategories
\[
 (J\downarrow\operatorname{Coalg}^{\mathrm{strong}})
 \simeq
 (\mathcal U_X\downarrow
       \operatorname{Alg}^{\mathrm{aug},Z}).
\]
Here the algebra on the right is augmented over \(K\) and its
augmentation ideal satisfies the support condition of
Theorem~\ref{thm:words-support}.  The two maps
\[
 \mathcal U_X\longrightarrow A,\qquad A\longrightarrow K
\]
have different meanings.  The first is a local vacuum marking;
the second is an empty-word augmentation.  Neither is a unital
scalar augmentation of a vertex algebra.

For example, the Heisenberg field with generator \(a\), vacuum
\(\mathbf1\), and nonzero level \(\kappa\) satisfies
\(a_{(1)}a=\kappa\mathbf1\).  A unital map to the scalar vertex
algebra would send the left side to zero and the right side to
\(\kappa\).  Such a map cannot exist.  The contradiction uses
preservation of the vertex product \(a_{(1)}a\), a different
condition from either geometric map.

\section{Finite summands and geometric extensions}

Suppose a coefficient is explicitly a finite direct sum indexed
by two decompositions of the same finite set.  The summands for
which those decompositions agree admit an inclusion and a
projection.  Both maps come from the given direct-sum presentation.

A collision stratum need not be a direct summand of a geometric
coefficient.  Already over \(R=k[u]\), the sequence
\[
 0\longrightarrow R\longrightarrow R[u^{-1}]
   \longrightarrow R[u^{-1}]/R\longrightarrow0
\]
does not split as \(R\)-modules.  If a retraction \(\rho\) existed,
then \(1=\rho(1)=u\rho(u^{-1})\), impossible in \(k[u]\).
Thus a finite indexing argument does not produce a retraction
across a collision divisor.  A geometric interchange requires
its actual maps, including any extension class supported on that
divisor.  Failure of this particular projection does not rule out
a different distributive law.

The cut product avoids choosing such a projection: its open
localizations and their base-change maps are specified in
\eqref{eq:words-cut-product}--\eqref{eq:words-diagonal-cut}.
Extending reconstruction from this fixed-curve word category to
arbitrary collision presentations still requires the comparisons
on those presentations.  The support calculation identifies what
they must retain, including crossed collisions and the local
vacuum.
